\chapter{Introdução}\label{cap:introducao}

\TextoModelo{Apresente o tema, delimite o problema geológico, sintetize o estado do conhecimento, justifique o estudo e antecipe a abordagem utilizada. A introdução deve conduzir naturalmente aos objetivos.}

Como exemplo de citação no sistema autor-data, uma referência pode ser integrada ao texto por meio de \citeonline{hasui2012} ou apresentada entre parênteses \cite{press2006}. Substitua estas referências pelas fontes efetivamente utilizadas.

\section{Objetivos}\label{sec:objetivos}

\subsection{Objetivo geral}

\TextoModelo{Declare o resultado científico principal que o trabalho pretende alcançar.}

\subsection{Objetivos específicos}

\begin{itemize}
  \item \TextoModelo{caracterizar as unidades geológicas da área de estudo;}
  \item \TextoModelo{descrever os procedimentos de campo, laboratório e gabinete;}
  \item \TextoModelo{integrar e interpretar os dados obtidos;}
  \item \TextoModelo{discutir as implicações geológicas dos resultados.}
\end{itemize}

\section{Justificativa}

\TextoModelo{Explique a relevância científica, regional, ambiental, econômica ou metodológica do estudo, sem repetir os objetivos.}
